\documentclass[aspectratio=169]{beamer}

\usetheme{Madrid}
\usecolortheme{default}
\usepackage[italian]{babel}
\usepackage[utf8]{inputenc}
\usepackage{listings}
\usepackage{xcolor}
\usepackage{booktabs}
\usepackage{array}
\usepackage{tikz}
\usetikzlibrary{arrows.meta, positioning, shapes.geometric}

% Colori personalizzati
\definecolor{sqlblue}{RGB}{0,70,127}
\definecolor{sqlgray}{RGB}{240,240,240}
\definecolor{codegreen}{RGB}{0,100,0}
\definecolor{codekeyword}{RGB}{0,0,180}
\definecolor{codegray}{RGB}{128,128,128}
\definecolor{codered}{RGB}{180,0,0}

% Stile listati SQL
\lstdefinestyle{sqlstyle}{
  language=SQL,
  basicstyle=\ttfamily\footnotesize,
  keywordstyle=\color{codekeyword}\bfseries,
  commentstyle=\color{codegreen}\itshape,
  stringstyle=\color{codered},
  backgroundcolor=\color{sqlgray},
  frame=single,
  framerule=0.5pt,
  rulecolor=\color{sqlblue!40},
  numbers=none,
  breaklines=true,
  showstringspaces=false,
  tabsize=2,
  morekeywords={CREATE,VIEW,AS,SELECT,FROM,WHERE,GROUP,BY,HAVING,COUNT,MAX,MIN,AVG,SUM,
                UPDATE,SET,INSERT,INTO,VALUES,WITH,UNION,ALL,TABLE,DISTINCT,JOIN,ON,
                LEFT,OUTER,CHECK,OPTION,CASCADED,LOCAL,AND,OR,NOT,IN,AS,RECURSIVE,LATERAL}
}

\lstset{style=sqlstyle}

% Impostazioni Beamer
\setbeamertemplate{navigation symbols}{}
\setbeamertemplate{footline}[frame number]
\setbeamerfont{frametitle}{size=\large, series=\bfseries}
\setbeamercolor{frametitle}{bg=sqlblue, fg=white}
\setbeamercolor{title}{fg=white}
\setbeamercolor{structure}{fg=sqlblue}
\setbeamercolor{block title}{bg=sqlblue!80, fg=white}
\setbeamercolor{block body}{bg=sqlblue!10}
\setbeamercolor{alerted text}{fg=codered}

\title[SQL: Viste e Table Expressions]{Il linguaggio SQL\\Viste e tabelle derivate}
\subtitle{Tutti i diritti riservati}
\author{Prof. Fedeli Massimo - Tutti i diritti riservati}
\date{2026}
\institute{Versione elettronica: 04.5.SQL.viste}

% -------------------------------------------------------
\begin{document}

% SLIDE 1 — Titolo
\begin{frame}
  \titlepage
\end{frame}

% -------------------------------------------------------
% SLIDE 2 — DB di riferimento
\begin{frame}{DB di riferimento per gli esempi}

% Tabella Imp — occupa tutta la larghezza
\textbf{Imp}\\[3pt]
{\small
\begin{tabular}{|l|l|l|l|r|}
  \hline
  \textbf{CodImp} & \textbf{Nome} & \textbf{Sede} & \textbf{Ruolo} & \textbf{Stipendio}\\
  \hline
  E001 & Rossi    & S01 & Analista      & 2000 \\
  E002 & Verdi    & S01 & Sistemista    & 1500 \\
  E003 & Bianchi  & S01 & Programmatore & 1000 \\
  E004 & Gialli   & S03 & Programmatore & 1000 \\
  E005 & Neri     & S02 & Analista      & 2500 \\
  E006 & Grigi    & S01 & Sistemista    & 1100 \\
  E007 & Violetti & S01 & Programmatore & 1000 \\
  E008 & Aranci   & S02 & Programmatore & 1200 \\
  \hline
\end{tabular}
}

\vspace{8pt}
% Sedi e Prog affiancate sotto
\begin{columns}[T]
\column{0.35\textwidth}
  \textbf{Sedi}\\[3pt]
  {\small
  \begin{tabular}{|l|l|}
    \hline
    \textbf{Sede} & \textbf{Città}\\
    \hline
    S01 & Bologna \\
    S02 & Milano  \\
    S03 & Roma    \\
    \hline
  \end{tabular}
  }
\column{0.55\textwidth}
  \textbf{Prog}\\[3pt]
  {\small
  \begin{tabular}{|l|l|l|}
    \hline
    \textbf{CodProg} & \textbf{Nome} & \textbf{Città}\\
    \hline
    P01 & Alpha & Bologna \\
    P02 & Beta  & Milano  \\
    \hline
  \end{tabular}
  }
\end{columns}

\end{frame}

% -------------------------------------------------------
% SLIDE 3 — Definizione di viste
\begin{frame}[fragile]{Definizione di viste}

\begin{block}{Concetto}
Mediante l'istruzione \texttt{CREATE VIEW} si definisce una ``tabella virtuale''.\\
Le tuple della vista sono il risultato di una query calcolata \alert{dinamicamente} ogni volta che si fa riferimento alla vista.
\end{block}

\vspace{6pt}
\begin{columns}[T]
\column{0.55\textwidth}
\begin{lstlisting}
CREATE VIEW ProgSedi(CodProg,CodSede)
AS SELECT P.CodProg, S.Sede
   FROM Prog P, Sedi S
   WHERE P.Citta = S.Citta

SELECT *
FROM ProgSedi
WHERE CodProg = 'P01'
\end{lstlisting}

\column{0.40\textwidth}
\vspace{8pt}
\textbf{Risultato:}\\[6pt]
\small
\begin{tabular}{|l|l|}
\hline
\textbf{CodProg} & \textbf{CodSede}\\
\hline
P01 & S01 \\
P01 & S03 \\
P01 & S02 \\
\hline
\end{tabular}
\end{columns}

\end{frame}

% -------------------------------------------------------
% SLIDE 4 — Uso delle viste
\begin{frame}{Uso delle viste}

Le viste possono essere create a vari scopi, tra cui i seguenti:

\begin{itemize}
  \item Permettere agli utenti di avere una \textbf{visione personalizzata} del DB, che in parte astragga dalla struttura logica del DB stesso
  \item Far fronte a \textbf{modifiche dello schema logico} che causerebbero la ricompilazione dei programmi applicativi
  \item \textbf{Semplificare la scrittura} di query complesse
\end{itemize}

\vspace{6pt}
\begin{block}{Nota}
Le viste possono essere usate come meccanismo di \textbf{controllo degli accessi}, fornendo ad ogni classe di utenti solo i privilegi necessari.\\
\vspace{4pt}
Nella definizione di una vista si possono usare anche \textbf{altre viste}.
\end{block}

\end{frame}

% -------------------------------------------------------
% SLIDE 5 — Indipendenza logica tramite viste
\begin{frame}[fragile]{Indipendenza logica tramite viste}

\begin{columns}[T]
\column{0.50\textwidth}
Si consideri un DB con:\\[4pt]
\texttt{EsamiSIT(Matr, Cognome, Nome, DataProva, Voto)}

\vspace{8pt}
Per evitare di ripetere i dati anagrafici, si decide di modificare lo schema, sostituendo alla tabella \texttt{EsamiSIT} due tabelle:
\begin{itemize}
  \item \texttt{StudentiSIT(Matr, Cognome, Nome)}
  \item \texttt{ProveSIT(Matr, DataProva, Voto)}
\end{itemize}

\column{0.48\textwidth}
È possibile ripristinare la ``visione originale'':
\begin{lstlisting}
CREATE VIEW EsamiSIT
  (Matr,Cognome,Nome,
   DataProva,Voto)
AS SELECT S.*, P.DataProva,
          P.Voto
   FROM StudentiSIT S,
        ProveSIT P
   WHERE S.Matr = P.Matr
\end{lstlisting}
\end{columns}

\end{frame}

% -------------------------------------------------------
% SLIDE 6 — Query complesse: esempio senza viste
\begin{frame}[fragile]{Query complesse che usano viste (1)}

\begin{block}{Esempio}
\textit{La sede che ha il massimo numero di impiegati}
\end{block}

\vspace{6pt}
\textbf{Soluzione senza viste:}

\begin{lstlisting}
SELECT I.Sede
FROM Imp I
GROUP BY I.Sede
HAVING COUNT(*) >= ALL
       (SELECT COUNT(*)
        FROM Imp I1
        GROUP BY I1.Sede)
\end{lstlisting}

\vspace{4pt}
$\Rightarrow$ La soluzione è corretta ma \alert{poco leggibile} e non intuitiva.

\end{frame}

% -------------------------------------------------------
% SLIDE 7 — Query complesse: soluzione con viste
\begin{frame}[fragile]{Query complesse che usano viste (2)}

\textbf{Soluzione con viste:}

\begin{lstlisting}
CREATE VIEW NumImp(Sede, NImp)
AS SELECT Sede, COUNT(*)
   FROM Imp
   GROUP BY Sede

SELECT Sede
FROM NumImp
WHERE NImp = (SELECT MAX(NImp)
              FROM NumImp)
\end{lstlisting}

\begin{block}{Osservazione}
Questa formulazione permette di trovare ``il MAX dei COUNT(*)'', operazione che \alert{non si può fare} direttamente scrivendo \texttt{MAX(COUNT(*))}.
\end{block}

\end{frame}

% -------------------------------------------------------
% SLIDE 8 — Query con più passi di raggruppamento
\begin{frame}[fragile]{Query complesse che usano viste (3)}

Con le viste è possibile risolvere query che richiedono \textbf{più passi di raggruppamento}:

\begin{block}{Esempio}
\textit{Per ogni valore (arrotondato) di stipendio medio per sede, il numero delle sedi che pagano tale stipendio medio}
\end{block}

\vspace{4pt}
Occorre aggregare prima per sede, poi per valore di stipendio medio:

\begin{lstlisting}
CREATE VIEW StipSedi(Sede, AvgStip)
AS SELECT Sede, AVG(Stipendio)
   FROM Imp
   GROUP BY Sede

SELECT AvgStip, COUNT(*) AS NumSedi
FROM StipSedi
GROUP BY AvgStip
\end{lstlisting}

\end{frame}

% -------------------------------------------------------
% SLIDE 9 — Aggiornamento di viste
\begin{frame}[fragile]{Aggiornamento di viste}

Le viste possono essere utilizzate per le interrogazioni come le normali tabelle del DB, ma per le \alert{operazioni di aggiornamento} sorgono problemi.

\begin{lstlisting}
CREATE VIEW NumImp(Sede, NImp)
AS SELECT Sede, COUNT(*)
   FROM Imp
   GROUP BY Sede

UPDATE NumImp
SET NImp = NImp + 1
WHERE Sede = 'S03'
\end{lstlisting}

\begin{alertblock}{Problema}
\textbf{Cosa significa incrementare il numero di impiegati di una sede?} \\
Non è possibile tradurre questa operazione in un aggiornamento della tabella \texttt{Imp}: \alert{non si può fare!}
\end{alertblock}

\end{frame}

% -------------------------------------------------------
% SLIDE 10 — Aggiornabilità di viste (1)
\begin{frame}{Aggiornabilità di viste (1)}

\begin{block}{Definizione formale}
Una vista è una funzione che calcola un risultato $r_V$ a partire da un'istanza di database $r$:
$$r_V = V(r)$$
L'aggiornamento di una vista (trasformare $r_V$ in $r'_V$) è possibile \textbf{solo se} è univocamente definita la nuova istanza $r'$ del DB — ovvero se la vista è \textbf{invertibile}.
\end{block}

\vspace{8pt}
Data la complessità del problema, ogni DBMS definisce quali viste sono aggiornabili.

\vspace{6pt}
\textbf{Restrizioni comuni} — il blocco più esterno della query di definizione non deve contenere:
\begin{itemize}
  \item \texttt{GROUP BY}
  \item Funzioni aggregate
  \item \texttt{DISTINCT}
  \item Join (espliciti o impliciti)
\end{itemize}

\end{frame}

% -------------------------------------------------------
% SLIDE 11 — Aggiornabilità di viste (2)
\begin{frame}[fragile]{Aggiornabilità di viste (2)}

La precisazione riguarda il \textbf{blocco più esterno} della query di definizione.

\vspace{6pt}
\textbf{Vista non aggiornabile} (join esplicito nel blocco esterno):
\begin{lstlisting}
CREATE VIEW ImpBO(CodImp,Nome,Sede,Ruolo,Stipendio)
AS SELECT I.*
   FROM Imp I JOIN Sedi S ON (I.Sede = S.Sede)
   WHERE S.Citta = 'Bologna'
\end{lstlisting}

\vspace{4pt}
\textbf{Vista aggiornabile} (join spostato in subquery):
\begin{lstlisting}
CREATE VIEW ImpBO(CodImp,Nome,Sede,Ruolo,Stipendio)
AS SELECT I.*
   FROM Imp I
   WHERE I.Sede IN (SELECT S.Sede
                    FROM Sedi S
                    WHERE S.Citta = 'Bologna')
\end{lstlisting}

\vspace{2pt}
Le due viste sono \textbf{logicamente equivalenti}, ma solo la seconda è aggiornabile.

\end{frame}

% -------------------------------------------------------
% SLIDE 12 — Viste con CHECK OPTION (1)
\begin{frame}[fragile]{Viste con CHECK OPTION (1)}

Per le viste aggiornabili si presenta un ulteriore problema. Si consideri:

\begin{lstlisting}
INSERT INTO ImpBO(CodImp,Nome,Sede,Ruolo,Stipendio)
VALUES ('E009','Azzurri','S03','Analista',2000)
\end{lstlisting}

\begin{alertblock}{Problema}
Il valore di \texttt{Sede} (\texttt{'S03'}) \alert{non rispetta} la specifica della vista \texttt{ImpBO} (solo impiegati di Bologna).
Una successiva query su \texttt{ImpBO} \textbf{non restituirebbe} la tupla appena inserita!
\end{alertblock}

\vspace{8pt}
\begin{block}{Soluzione}
All'atto della definizione della vista si può specificare la clausola \texttt{WITH CHECK OPTION}, che impone che ogni tupla inserita nella vista sia anche \textbf{restituita dalla sua query di definizione}.
\end{block}

\end{frame}

% -------------------------------------------------------
% SLIDE 13 — Viste con CHECK OPTION (2)
\begin{frame}[fragile]{Viste con CHECK OPTION (2)}

\begin{lstlisting}
CREATE VIEW ImpBO(CodImp,Nome,Sede,Ruolo,Stipendio)
AS SELECT I.*
   FROM Imp I
   WHERE I.Sede IN (SELECT S.Sede
                    FROM Sedi S
                    WHERE S.Citta = 'Bologna')
WITH CHECK OPTION
\end{lstlisting}

\vspace{10pt}
In questo modo l'inserimento:

\begin{lstlisting}
INSERT INTO ImpBO(CodImp,Nome,Sede,Ruolo,Stipendio)
VALUES ('E009','Azzurri','S03','Analista',2000)
\end{lstlisting}

\begin{alertblock}{}
\centering viene \textbf{impedito} dal DBMS
\end{alertblock}

\end{frame}

% -------------------------------------------------------
% SLIDE 14 — Tipi di CHECK OPTION (1)
\begin{frame}{Tipi di CHECK OPTION (1)}

Si considera il caso in cui una vista \textbf{V1} sia definita in termini di un'altra vista \textbf{V2}.

\vspace{8pt}
\begin{block}{\texttt{WITH CASCADED CHECK OPTION} (default)}
Se si crea V1 con \texttt{WITH CHECK OPTION}, il DBMS verifica che la nuova tupla $t$ soddisfi sia la specifica di V1 \textbf{sia quella di V2} (e di tutte le altre viste dipendenti), indipendentemente dal fatto che V2 sia stata definita con \texttt{CHECK OPTION}.\\
\vspace{4pt}
È il comportamento di default: equivalente a \texttt{WITH CASCADED CHECK OPTION}.
\end{block}

\vspace{6pt}
\begin{block}{\texttt{WITH LOCAL CHECK OPTION}}
Il DBMS verifica solo che $t$ soddisfi la specifica di V1 e delle sole viste da cui V1 dipende per cui è stata \textbf{esplicitamente} specificata la \texttt{CHECK OPTION}.
\end{block}

\end{frame}

% -------------------------------------------------------
% SLIDE 15 — Tipi di CHECK OPTION (2)
\begin{frame}{Tipi di CHECK OPTION (2)}

\begin{columns}[T]
\column{0.45\textwidth}
\textbf{Grafo di dipendenze tra viste:}

\vspace{12pt}
\begin{tikzpicture}[
  node distance=1.2cm,
  box/.style={draw, rectangle, rounded corners=3pt, minimum width=1.4cm, minimum height=0.7cm, fill=sqlblue!15}
]
\node[box] (R) {R};
\node[box, below left=of R] (V1) {V1};
\node[box, below right=of R] (V2) {V2};
\node[box, below=of V1] (V3) {V3};
\node[box, below=of V2] (V4) {V4};

\draw[-{Stealth}] (R) -- (V1);
\draw[-{Stealth}] (R) -- (V2);
\draw[-{Stealth}] (V1) -- (V3);
\draw[-{Stealth}] (V2) -- (V4);
\draw[-{Stealth}] (V3) -- (V4);
\end{tikzpicture}

\column{0.52\textwidth}
\textbf{Viste controllate all'inserimento:}\\[8pt]
\small
\begin{tabular}{|l|l|}
\hline
\textbf{INSERT INTO} & \textbf{Si controllano\ldots}\\
\hline
V1 & -- \\
V2 & -- \\
V3 & V1, V3 \\
V4 & V1, V2, V4 \\
\hline
\end{tabular}

\vspace{10pt}
\footnotesize
(Assumendo: V3 = LOCAL, V4 = CASCADED, V1 e V2 senza CHECK OPTION)
\end{columns}

\end{frame}

% -------------------------------------------------------
% SLIDE 16 — Table expressions (1)
\begin{frame}[fragile]{Table expressions (1)}

Una delle caratteristiche più interessanti di SQL è la possibilità di usare all'interno della clausola \texttt{FROM} una \textbf{subquery} che definisce dinamicamente una tabella derivata: \textbf{table expression}.

\begin{block}{Esempio}
\textit{Per ogni sede, lo stipendio massimo e quanti impiegati lo percepiscono}
\end{block}

\begin{lstlisting}
SELECT SM.Sede, SM.MaxStip, COUNT(*) AS N
FROM Imp I,
     (SELECT Sede, MAX(Stipendio) AS MaxStip
      FROM Imp
      GROUP BY Sede) AS SM(Sede, MaxStip)
WHERE I.Sede = SM.Sede
  AND I.Stipendio = SM.MaxStip
GROUP BY SM.Sede, SM.MaxStip
\end{lstlisting}

\end{frame}

% -------------------------------------------------------
% SLIDE 17 — Table expressions (2)
\begin{frame}[fragile]{Table expressions (2)}

Le table expressions possono essere usate per realizzare \textbf{diversi passi di aggregazione}:

\begin{block}{Esempio}
\textit{Per ogni valore (arrotondato) di stipendio medio per sede, il numero delle sedi che pagano tale stipendio medio}
\end{block}

\begin{lstlisting}
SELECT AvgStip, COUNT(*) AS NumSedi
FROM (SELECT Sede, AVG(Stipendio)
      FROM Imp
      GROUP BY Sede) AS StipSedi(Sede, AvgStip)
GROUP BY AvgStip
\end{lstlisting}

\vspace{4pt}
$\Rightarrow$ Equivalente alla soluzione con le viste, ma senza definire una vista persistente.

\end{frame}

% -------------------------------------------------------
% SLIDE 18 — Table expressions correlate (1)
\begin{frame}[fragile]{Table expressions correlate (1)}

Una table expression può essere \textbf{correlata} a un'altra tabella presente nella clausola \texttt{FROM}. In DB2 si usa la parola riservata \texttt{TABLE}.

\begin{block}{Esempio}
\textit{Per ogni sede, la somma degli stipendi pagati agli analisti}
\end{block}

\begin{lstlisting}
SELECT S.Sede, Stip.TotStip
FROM Sedi S,
     TABLE(SELECT SUM(Stipendio)
           FROM Imp I
           WHERE I.Sede = S.Sede
             AND I.Ruolo = 'Analista') AS Stip(TotStip)
\end{lstlisting}

\vspace{4pt}
\begin{block}{Nota}
Le sedi senza analisti \textbf{compaiono} nell'output con \texttt{TotStip = NULL}. Usando il \texttt{GROUP BY} lo stesso risultato richiederebbe un \texttt{LEFT OUTER JOIN}, con attenzione ai valori nulli.
\end{block}

\end{frame}

% -------------------------------------------------------
% SLIDE 19 — Table expressions correlate (2)
\begin{frame}[fragile]{Table expressions correlate (2)}

\begin{block}{Esempio}
\textit{Per ogni sede, il numero di analisti e la somma degli stipendi}
\end{block}

\begin{lstlisting}
SELECT S.Sede, Stip.NumAn, Stip.TotStip
FROM Sedi S,
     TABLE(SELECT COUNT(*), SUM(Stipendio)
           FROM Imp I
           WHERE I.Sede = S.Sede
             AND I.Ruolo = 'Analista')
     AS Stip(NumAn, TotStip)
\end{lstlisting}

\vspace{4pt}
Per le sedi senza analisti: \texttt{NumAn = 0}, \texttt{TotStip = NULL}.

\vspace{4pt}
Con \texttt{LEFT OUTER JOIN} e \texttt{GROUP BY}, per le sedi senza analisti \texttt{TotStip} è nullo, ma \texttt{COUNT(*)} darebbe \textbf{1} (non 0), perché c'è una tupla nel risultato con valori nulli. Occorre usare \texttt{COUNT(CodImp)}.

\end{frame}

% -------------------------------------------------------
% SLIDE 20 — Limiti delle table expressions
\begin{frame}[fragile]{Limiti delle table expressions}

\begin{block}{Esempio}
\textit{La sede in cui la somma degli stipendi è massima}
\end{block}

\begin{lstlisting}
SELECT Sede
FROM (SELECT Sede, SUM(Stipendio) AS TotStip
      FROM Imp
      GROUP BY Sede) AS SediStip
WHERE TotStip = (SELECT MAX(TotStip)
                 FROM (SELECT Sede, SUM(Stipendio) AS TotStip
                       FROM Imp
                       GROUP BY Sede) AS SediStip2)
\end{lstlisting}

\begin{alertblock}{Problema}
Benché la query sia corretta, le due table expressions sono \textbf{identiche}: la stessa subquery viene calcolata \alert{due volte}, portando a inefficienza e scarsa leggibilità.
\end{alertblock}

\end{frame}

% -------------------------------------------------------
% SLIDE 21 — Common table expressions
\begin{frame}[fragile]{Common table expressions}

L'idea alla base delle \textbf{common table expressions} (CTE) è definire una ``vista temporanea'' che può essere usata in una query come se fosse una \texttt{VIEW}.

\begin{lstlisting}
WITH SediStip(Sede, TotStip)
AS (SELECT Sede, SUM(Stipendio)
    FROM Imp
    GROUP BY Sede)

SELECT Sede
FROM SediStip
WHERE TotStip = (SELECT MAX(TotStip)
                 FROM SediStip)
\end{lstlisting}

\begin{block}{Sintassi generale}
La clausola \texttt{WITH} supporta la definizione di \textbf{più CTE}:
\begin{lstlisting}
WITH CTE1(...) AS (...),
     CTE2(...) AS (...)
SELECT ...
\end{lstlisting}
\end{block}

\end{frame}

% -------------------------------------------------------
% SLIDE 22 — WITH e interrogazioni ricorsive (1)
\begin{frame}[fragile]{WITH e interrogazioni ricorsive (1)}

\begin{columns}[T]
\column{0.55\textwidth}
Si consideri la tabella \texttt{Genitori(Figlio, Genitore)}.

\begin{block}{Esempio}
\textit{Trova tutti gli antenati (genitori, nonni, bisnonni, \ldots) di una persona}
\end{block}

La query è \textbf{ricorsiva}: richiede un numero di (self-)join non noto a priori — non è esprimibile con SQL standard.

\vspace{6pt}
\textbf{La CTE ricorsiva} \texttt{Antenati(Persona, Avo)} è composta da:
\begin{itemize}
  \item Una \textbf{subquery base} (non ricorsiva) che inizializza con i dati di \texttt{Genitori}
  \item Una \textbf{subquery ricorsiva} che aggiunge le coppie risultanti dal join tra \texttt{Genitori} e \texttt{Antenati}
\end{itemize}

\column{0.42\textwidth}
\textbf{Genitori:}\\[4pt]
\small
\begin{tabular}{|l|l|}
\hline
\textbf{Figlio} & \textbf{Genitore}\\
\hline
Anna    & Luca   \\
Luca    & Maria  \\
Luca    & Giorgio\\
Giorgio & Lucia  \\
\hline
\end{tabular}

\vspace{10pt}
\textbf{Antenati (risultato):}\\[4pt]
\small
\begin{tabular}{|l|l|}
\hline
\textbf{Persona} & \textbf{Avo}\\
\hline
Anna & Luca    \\
Luca & Maria   \\
Luca & Giorgio \\
Giorgio & Lucia \\
Anna & Maria   \\
Anna & Giorgio \\
Anna & Lucia   \\
Luca & Lucia   \\
\hline
\end{tabular}
\end{columns}

\end{frame}

% -------------------------------------------------------
% SLIDE 23 — WITH e interrogazioni ricorsive (2)
\begin{frame}[fragile]{WITH e interrogazioni ricorsive (2)}

\begin{lstlisting}
WITH Antenati(Persona, Avo)
AS (
  (SELECT Figlio, Genitore   -- subquery base
   FROM Genitori)
  UNION ALL
  (SELECT G.Figlio, A.Avo   -- subquery ricorsiva
   FROM Genitori G, Antenati A
   WHERE G.Genitore = A.Persona)
)

SELECT Avo
FROM Antenati
WHERE Persona = 'Anna'
\end{lstlisting}

\begin{block}{Nota}
L'uso di \texttt{UNION ALL} garantisce che tutte le tuple vengano propagate ad ogni iterazione, comprese le eventuali duplicate.
\end{block}

\end{frame}

% -------------------------------------------------------
% SLIDE 24 — Valutazione della ricorsione
\begin{frame}[fragile]{WITH e interrogazioni ricorsive (3) — Valutazione}

Ad ogni iterazione il DBMS aggiunge ad \texttt{Antenati} le tuple risultanti dal join tra \texttt{Genitori} e le \textbf{sole tuple aggiunte} all'iterazione precedente (semantica ``delta'').

\vspace{8pt}
\begin{columns}[T]
\column{0.42\textwidth}
\textbf{Iterazione 0 (base):}\\[4pt]
\small
\begin{tabular}{|l|l|}
\hline
Anna & Luca \\ Luca & Maria \\ Luca & Giorgio \\ Giorgio & Lucia \\
\hline
\end{tabular}

\vspace{8pt}
\textbf{Iterazione 1 (ricorsiva):}\\[4pt]
\small
\begin{tabular}{|l|l|}
\hline
Anna & Maria \\ Anna & Giorgio \\ Luca & Lucia \\
\hline
\end{tabular}

\column{0.42\textwidth}
\textbf{Iterazione 2 (ricorsiva):}\\[4pt]
\small
\begin{tabular}{|l|l|}
\hline
Anna & Lucia \\
\hline
\end{tabular}

\vspace{8pt}
\textbf{Iterazione 3:}\\[4pt]
\small Nessuna nuova tupla $\Rightarrow$ \textbf{STOP}

\vspace{8pt}
\begin{block}{}
La terminazione è garantita se il DB è \textbf{aciclico} (DAG).
\end{block}
\end{columns}

\end{frame}

% -------------------------------------------------------
% SLIDE 25 — Tipi notevoli di query ricorsive
\begin{frame}{Tipi notevoli di query ricorsive}

\begin{block}{Caso base (già visto)}
DB aciclico (DAG): la terminazione è garantita e non si calcola nessuna informazione sui ``percorsi'' trovati.
\end{block}

\vspace{8pt}
\textbf{Percorso}: sequenza di tuple del DB che generano una tupla nel risultato finale (nell'esempio, una specifica ``linea genealogica'').

\vspace{12pt}
\textbf{Due casi notevoli da esaminare:}

\begin{enumerate}
  \item \textbf{Query ricorsive con informazione sui percorsi} (lunghezza, costo cumulativo, \ldots)
  \item \textbf{Condizione di stop} per query ricorsive su DB con \textbf{cicli}
\end{enumerate}

\end{frame}

% -------------------------------------------------------
% SLIDE 26 — Informazione sui percorsi (1)
\begin{frame}{Informazione sui percorsi (1)}

Il caso più semplice prevede l'aggiunta di un'informazione sulla \textbf{lunghezza} (distanza, livello, ecc.) del percorso.

\vspace{8pt}
\begin{columns}[T]
\column{0.42\textwidth}
\textbf{Genitori:}\\[4pt]
\small
\begin{tabular}{|l|l|}
\hline
\textbf{Figlio} & \textbf{Genitore}\\
\hline
Anna & Luca \\ Luca & Maria \\ Luca & Giorgio \\ Giorgio & Lucia \\
\hline
\end{tabular}

\column{0.52\textwidth}
\textbf{Antenati (con lunghezza):}\\[4pt]
\small
\begin{tabular}{|l|l|l|}
\hline
\textbf{Persona} & \textbf{Avo} & \textbf{Lungh}\\
\hline
Anna & Luca    & 1 \\
Luca & Maria   & 1 \\
Luca & Giorgio & 1 \\
Giorgio & Lucia & 1 \\
Anna & Maria   & 2 \\
Anna & Giorgio & 2 \\
Luca & Lucia   & 2 \\
Anna & Lucia   & 3 \\
\hline
\end{tabular}
\end{columns}

\end{frame}

% -------------------------------------------------------
% SLIDE 27 — Informazione sui percorsi (2)
\begin{frame}[fragile]{Informazione sui percorsi (2)}

La lunghezza parte da \textbf{1} e si incrementa di 1 ad ogni passo ricorsivo:

\begin{lstlisting}
WITH Antenati(Persona, Avo, Lungh)
AS (
  (SELECT Figlio, Genitore, 1  -- caso base
   FROM Genitori)
  UNION ALL
  (SELECT G.Figlio, A.Avo, A.Lungh + 1
   FROM Genitori G, Antenati A
   WHERE G.Genitore = A.Persona)
)

SELECT *
FROM Antenati
WHERE Persona = 'Anna'
\end{lstlisting}

\end{frame}

% -------------------------------------------------------
% SLIDE 28 — Informazione sui percorsi (3) — BOM
\begin{frame}{Informazione sui percorsi (3) — Bill of Materials}

In alcuni casi l'informazione sui percorsi è un \textbf{costo cumulativo moltiplicativo}:

\begin{block}{Esempio: relazione PARTI}
Per ogni parte si conoscono le sotto-parti e le quantità necessarie per assemblarla.\\
\textit{Si vuole sapere quante unità di una parte servono in totale per costruire un'altra parte.}
\end{block}

\vspace{8pt}
\begin{columns}[T]
\column{0.45\textwidth}
\textbf{Parti:}\\[4pt]
\small
\begin{tabular}{|l|l|l|}
\hline
\textbf{Parte} & \textbf{Subparte} & \textbf{Qta}\\
\hline
P1 & P2 & 2 \\
P1 & P3 & 3 \\
P2 & P4 & 5 \\
P3 & P4 & 4 \\
\hline
\end{tabular}

\column{0.50\textwidth}
\vspace{4pt}
Per costruire P1 occorrono:
\begin{itemize}
  \item Percorso P1--P2--P4: $2 \times 5 = 10$ unità di P4
  \item Percorso P1--P3--P4: $3 \times 4 = 12$ unità di P4
  \item Totale P4: $10 + 12 = \mathbf{22}$
\end{itemize}
\end{columns}

\end{frame}

% -------------------------------------------------------
% SLIDE 29 — Informazione sui percorsi (4) — BOM schema
\begin{frame}{Informazione sui percorsi (4) — Schema soluzione}

La soluzione prevede due fasi:

\begin{enumerate}
  \item \textbf{Ricorsione} per calcolare l'informazione di ogni singolo percorso (prodotto delle quantità):
    \begin{itemize}
      \item P1--P2--P4: $2 \times 5 = 10$
      \item P1--P3--P4: $3 \times 4 = 12$
    \end{itemize}
  \item \textbf{Aggregazione} dei risultati dei singoli percorsi sulla vista ricorsiva (\texttt{GROUP BY + SUM})
\end{enumerate}

\vspace{8pt}
\begin{columns}[T]
\column{0.48\textwidth}
\textbf{Percorsi (dopo ricorsione):}\\[4pt]
\small
\begin{tabular}{|l|l|l|}
\hline
\textbf{Comp.} & \textbf{Comp.} & \textbf{Qty}\\
\hline
P1 & P2 & 2 \\
P1 & P3 & 3 \\
P2 & P4 & 5 \\
P3 & P4 & 4 \\
P1 & P4 & 10 \\
P1 & P4 & 12 \\
\hline
\end{tabular}
\column{0.48\textwidth}
\textbf{Risultato finale (dopo GROUP BY):}\\[4pt]
\small
\begin{tabular}{|l|l|l|}
\hline
\textbf{Comp.} & \textbf{Comp.} & \textbf{Qt. Totale}\\
\hline
P1 & P2 & 2  \\
P1 & P3 & 3  \\
P2 & P4 & 5  \\
P3 & P4 & 4  \\
P1 & P4 & 22 \\
\hline
\end{tabular}
\end{columns}

\end{frame}

% -------------------------------------------------------
% SLIDE 30 — Informazione sui percorsi (5) — Query BOM
\begin{frame}[fragile]{Informazione sui percorsi (5) — Query BOM}

Ad ogni passo si \textbf{moltiplica} la quantità per quella del nodo aggiunto:

\begin{lstlisting}
WITH Percorsi(Composto, Componente, Qty)
AS (
  (SELECT Parte, Subparte, Qta
   FROM Parti)
  UNION ALL
  (SELECT H.Composto, P.Subparte, H.Qty * P.Qta
   FROM Parti P, Percorsi H
   WHERE H.Componente = P.Parte)
)

SELECT Composto, Componente, SUM(Qty) AS QtTotale
FROM Percorsi
GROUP BY Composto, Componente
\end{lstlisting}

\end{frame}

% -------------------------------------------------------
% SLIDE 31 — Database ciclici
\begin{frame}{Database ciclici}

Nel caso di DB \textbf{ciclici} occorre prestare particolare attenzione alla terminazione.

\begin{alertblock}{Rischio}
Senza condizioni di stop, la ricorsione può non terminare mai!
\end{alertblock}

\vspace{8pt}
Non disponendo di altri strumenti, occorre inserire nella subquery ricorsiva un \textbf{predicato} che sia falso per tutte le nuove tuple che si andrebbero ad aggiungere, fermando così l'iterazione.

\vspace{8pt}
\textbf{Casi tipici di condizione di stop:}
\begin{itemize}
  \item \textbf{Limitazione sulla lunghezza massima} dei percorsi
  \item \textbf{Limitazione sul ``costo''} (somma dei pesi degli archi)
\end{itemize}

\end{frame}

% -------------------------------------------------------
% SLIDE 32 — Limitazione sulla lunghezza (1)
\begin{frame}{Limitazione sulla lunghezza dei percorsi (1)}

\begin{columns}[T]
\column{0.55\textwidth}
Il grafo contiene un \textbf{ciclo} (ad esempio: A--B--E--A--B--E--A--\ldots).

\vspace{8pt}
Fissando un \textbf{limite sulla lunghezza} dei percorsi si garantisce che ogni percorso trovato abbia lunghezza finita, quindi l'esecuzione prima o poi termina.

\vspace{8pt}
\begin{block}{Attenzione}
Il limite va scelto in modo da includere tutti i percorsi ``interessanti''.
\end{block}

\column{0.40\textwidth}
\vspace{10pt}
\begin{tikzpicture}[
  node distance=1.3cm,
  nd/.style={draw, circle, fill=sqlblue!20, minimum size=0.7cm}
]
\node[nd] (A) {A};
\node[nd, below right=of A] (B) {B};
\node[nd, right=of A] (C) {C};
\node[nd, below right=of B] (D) {D};
\node[nd, below left=of B] (E) {E};

\draw[-{Stealth}] (A) -- (B);
\draw[-{Stealth}] (A) -- (C);
\draw[-{Stealth}] (B) -- (D);
\draw[-{Stealth}] (B) -- (E);
\draw[-{Stealth}] (E) -- (A);
\end{tikzpicture}
\end{columns}

\end{frame}

% -------------------------------------------------------
% SLIDE 33 — Limitazione sulla lunghezza (2)
\begin{frame}[fragile]{Limitazione sulla lunghezza dei percorsi (2)}

Con limite a \textbf{3} (e condizione anti-ciclo \texttt{P.Da <> S.To}):

\begin{lstlisting}
WITH Percorsi(Da, A, Lungh)
AS (
  (SELECT From, To, 1
   FROM Siti)
  UNION ALL
  (SELECT P.Da, S.To, P.Lungh + 1
   FROM Siti S, Percorsi P
   WHERE P.A = S.From
     AND P.Da <> S.To
     AND P.Lungh + 1 <= 3)
)

SELECT *
FROM Percorsi
\end{lstlisting}

\vspace{4pt}
La condizione \texttt{P.Da <> S.To} serve ad evitare di includere un percorso \textbf{inutile} nel risultato (ciclo che torna al punto di partenza).

\end{frame}

% -------------------------------------------------------
% SLIDE 34 — Limitazione sulla lunghezza (3) — Nota
\begin{frame}{Limitazione sulla lunghezza dei percorsi (3) — Nota}

\begin{columns}[T]
\column{0.55\textwidth}
La condizione \texttt{P.Da <> S.To} serve ad evitare di includere percorsi inutili nel risultato.

\vspace{8pt}
\begin{alertblock}{Attenzione}
Benché in alcuni casi particolari possa essere usata come condizione di terminazione, \textbf{nel caso generale non è sufficiente} a garantire la terminazione in DB ciclici.
\end{alertblock}

\vspace{8pt}
Il grafo di esempio mostra un caso in cui il ciclo è rilevabile, ma cicli più lunghi o complessi richiedono comunque il limite esplicito sulla lunghezza.

\column{0.40\textwidth}
\textbf{Siti:}\\[6pt]
\small
\begin{tabular}{|l|l|}
\hline
\textbf{From} & \textbf{To}\\
\hline
A & B \\
B & C \\
B & C \\
C & B \\
\hline
\end{tabular}

\vspace{10pt}
\begin{tikzpicture}[
  node distance=1.1cm,
  nd/.style={draw, circle, fill=sqlblue!20, minimum size=0.6cm}
]
\node[nd] (A) {A};
\node[nd, right=of A] (B) {B};
\node[nd, right=of B] (C) {C};
\draw[-{Stealth}] (A) -- (B);
\draw[-{Stealth}, bend left=20] (B) to (C);
\draw[-{Stealth}, bend left=20] (C) to (B);
\end{tikzpicture}
\end{columns}

\end{frame}

% -------------------------------------------------------
% SLIDE 35 — Limitazione sul costo
\begin{frame}[fragile]{Limitazione sul ``costo''}

In maniera analoga si ragiona se si vuole imporre un \textbf{costo massimo} sui percorsi:

\begin{lstlisting}
WITH Percorsi(Da, A, TotKm)
AS (
  (SELECT From, To, Km
   FROM Paesi)
  UNION ALL
  (SELECT P.Da, S.To, P.TotKm + S.Km
   FROM Paesi S, Percorsi P
   WHERE P.A = S.From
     AND P.Da <> S.To
     AND P.TotKm + S.Km <= 17)
)

SELECT *
FROM Percorsi
\end{lstlisting}

\begin{block}{Osservazione}
In generale ci possono essere \textbf{più percorsi} tra 2 siti; va considerato quello a costo minore (problema del \textit{cammino minimo}).
\end{block}

\end{frame}

% -------------------------------------------------------
% SLIDE 36 — Riepilogo
\begin{frame}{Riepilogo}

\begin{block}{Viste (\texttt{CREATE VIEW})}
Tabelle virtuali, interrogabili come le tabelle reali. Soggette a limiti per quanto riguarda gli \textbf{aggiornamenti}.
\end{block}

\begin{block}{Table expressions}
Subquery che definisce una tabella derivata, utilizzabile nella clausola \texttt{FROM} come una normale tabella.
\end{block}

\begin{block}{Common table expressions (\texttt{WITH})}
``Vista temporanea'' usata in una query come se fosse a tutti gli effetti una \texttt{VIEW}. Evita la ridondanza delle table expressions identiche.
\end{block}

\begin{block}{Query ricorsive (\texttt{WITH RECURSIVE})}
Mediante CTE è possibile esprimere interrogazioni ricorsive, definendo ``viste temporanee ricorsive'' come unione di una \textbf{subquery base} e una \textbf{subquery ricorsiva}. Richiedono attenzione alla terminazione in DB ciclici.
\end{block}

\end{frame}

\end{document}
